\section{Conclusion}
\label{sec:conclusion}

The most useful thing an audit can do is put a floor under numbers people are standing behind. Computing $\mathrm{score}_5$ for a token bag and an untrained encoder on the 14 SemVec corpora, under those authors' own procedure, establishes a non-learned baseline floor: nothing we built reaches EqNet or the headline system on any corpus. We report that having first retracted the opposite conclusion, which our own non-matching protocol had produced.

A per-symbol tokenization convention makes physics-derived held-out-form benchmarks solvable by lookup ($84\%$ unique variable sets on AI~Feynman): design hygiene, not a field-level indictment, since the published equivalence-class family defines classes over a shared variable pool and is immune at tier~1. \textbf{The generalizable finding is tier~2}: operator-bag separability runs $66$--$100\%$ across every family we audited, including corpora perfectly clean at tier~1. Unlike tier~1 this is meaningful information, not a cheat: the problem is the mismatch between cue and claim. Removing variable leakage does not fix the evaluation.

Two positive results give the account its shape. First, sharing variables across classes removes the tier-1 cue, and a triplet encoder then identifies held-out forms at $0.932$ against $0.604$ for a random encoder and $0.56$ for a variable-blind bag (\S\ref{sec:positive_control}). Second, on EQNET \texttt{poly8} (1102 classes with a shared variable pool and zero unique operator bags), a Tree-LSTM reaches unseen-class $k$-NN precision of $0.896$ at $K{=}500$ against the strongest non-learned baseline's $0.518$, with the margin growing as $K$ grows, and reaches $0.994$ on a twin probe where token inventory \emph{and} local tree shape are both pinned at chance by construction (\S\ref{sec:structural_scale}). This clears the criterion of \S\ref{sec:tiers}: the highest floor we can build, not a demonstration of composition (see below). The bound is honest rather than decomposed: an untrained recursive encoder already reaches $0.788$ on that probe, so training is credited with at most the remaining $0.206$.

We applied the same standard to ourselves, and it cost us two claims: both purpose-built overlap suites turned out to be solvable at tier~2 (Appendix~P), and the protocol-sensitivity contrast shrank from $0.44$ to $0.13$ under corrected provenance (Table~\ref{tab:diversity_sweep}). Coverage is a causal part of what remains (\S\ref{sec:loso}).

\paragraph{What we claim, exactly.} The hierarchy is a \emph{falsification} device rather than an identification one: failing a tier invalidates the stronger reading, while clearing one establishes only insensitivity to the cues actually tested. \textbf{Established}: physics-derived held-out-form benchmarks admit a complete lower-tier shortcut. \textbf{Established}: the trained encoder exceeds every non-learned baseline we could construct, including one that holds inventory, local shape \emph{and} order fixed by construction. \textbf{Not established}: that the encoder performs compositional algebraic reasoning; no probe separates composition from a structural statistic we did not think to measure, since a sufficiently rich set of local relations determines the tree. A claim is worth exactly the highest floor it clears.

\paragraph{Limitations.} Classes have ${\approx}3$ members differing by one rewrite; claims concern equations, not seeds; the cautionary claim is scoped to physics corpora under per-symbol tokenization. Code, auditor, run tags, and a script that recomputes and asserts every non-learned number are included in the supplementary material.
